% For the drought impact base the research outputs, mainly consist of an spatial temparal nalysis of the data.
% The goal of this chapter is to provide an birds eye vieuw of all the data we have collected. What biases are in our data and what pattern's do frequenly occur in our data? Are there any temporal biases, are there any spatial biases. If we would build an model in the future, where should we start, what can we infer based of our simple plots. 



% Firstly in this section we provide an bird-eye vieuw of the netlherands. We plot some general stats about our database. Spatial, and temporal.  insert plots >spatial aggreated counts, what regions have the most impacts<, 
% and temporal how do impcats counts in genral evolve over time.


% Then we coonclude like that we have to do an data split based of time, and we sort of identify , two reporting regions pre-2018, and after 2018. Based on these splits, we will try to check how they correlate compared to the kmidrought composlite index, with an nrmolized impact ocunts plot for both feature's here we mianly ask, do the impact counts overlapp with drought periods? 




% In chpater valiation, we validated the data accros other proxy datasets. 

% The goal of this section, is different. It sort of has to provide an births-eye view of all the data we collected. The main question this should anwser to our stakeholder's is: what could we potentially do with this dataset? Do we have enough impacts to build an machine learning model? What drought impacts occur where and how often and at the right timing? 


% First of we start by analysing the general drought impact trend from our news-archive's. 

% >Insert plot<
% Heat-map per category + average impacts per year +

% years, drought index, category normalized.
% ------------------------------







% to find out how many impact we have and if they are good enough for machine learning, we require our impacts to be spatially and temporal be distrubted in an logical and sound way. 

% 1. total impact counts per category
% 2. mean variance per year














% Observe that there is an relationship between how much reporting there is done in general in the Netherlands and the amount of impact counts. The discrepancy between reported impacts in 2005 and 2020, could also be explained by different reporting practise's (the observered world) and not an actual changes in the real world.


% We therefor, present the data in 



% So we also want to visualize general trends in the impact database, but to do so we require





% then 



% First of how much data do we have in general, how much data do we have for a certain category. This gives us insights, in the spatial properties of the 

% Do general trends, in the data set compared to drought events make sense.

% Decide what visulizations we need:


% 0. General news-artilce counts 

% ---> motivate our impact aggregatoin method


% 1. amount of impacts for an specicif sub-category over year.

% Do normalized theme, and at an drougt index? maybe min max normalize for +/- 2.5 years,

% idea is to have an local and global variable. 



% 2. Percentage impact sub-plot.

% What regions have highest impact occurence's. so bassilcy we +/- for an specif impact an NL plot.

% these should be 20 plots. ---> use nuts3 aggregation framework to excute this.

% ---------------------------------------------------
% tldr.
